\paragraph{Conjugate-coordinate class: the next-to-leading symbol of the mass-supertrace term at order \(\epsilon^2\) (conditional; bounded no-go with a conditional Einstein--Hilbert-form limit).}
Keep the setting of the preceding paragraph and the term \(+M_{\mathrm{str}}\operatorname{Tr}[\chi\,(H_X-\mathrm i0)^{-1}]\), \(H_X=-\Box_X\). Take \(A^a=\epsilon\alpha^a\cos\phi\), \(\phi=q\cdot y\), \(W_a\) the Hamiltonian field of \(T^a\), \(Y=q_\mu\{y^\mu,\cdot\}\), \(\kappa_a=W_a(\phi)\) (so \(Y(\phi)=0\)) and the Gaussian window of \(\Delta_{3/4}\). Terms of fourth order in \(q\) at order \(\epsilon^2\) contain all derivatives on \(A\); they are computed with the first jet of \(T\) frozen and the window symbol constant, since window gradients and \(\partial^2T\) carry fewer powers of \(q\) (derivative counting; the window symbol being constant makes the composition terms with \(\chi\) total derivatives). In this sector \(H_X=\eta^{ab}\Pi_a\Pi_b\), \(\Pi_a=\omega_a-\epsilon\alpha_a\zeta\sin\phi\), \(\omega_a=\xi(W_a)\), \(\zeta=\xi(Y)\), \([\omega_a,\phi]=-\mathrm i\kappa_a\), with \(\zeta\) commuting with everything.

\emph{Symbol and checks.} The weight-two symbol of the resolvent was derived by the left-symbol recursion, expanded to \(\epsilon^2\) and averaged over \(\phi\) (the weight-one terms average to zero). It reproduces (a) the one-dimensional coefficient \(-\tfrac1{12}V'^2\) and the two-dimensional Maxwell coefficient \(-\tfrac1{12}F_{ab}F^{ab}\) to \(10^{-11}\); (b) the leading-symbol Hessian of the preceding paragraph; (c) in the non-degenerate case, with the vector fields divergence free, the Seeley--DeWitt coefficient \(\sqrt{|g|}(R/6-\tfrac14(\nabla\lambda)^2)\), \(\lambda=\ln\sqrt{|g|}\), at second order (Fierz--Pauli form), with relative phase \(\mathrm i\) to \(12\) digits.

\emph{Result in the class.} The integrand \(K(s)\) of the proper time \(s\) is finite at small \(s\) and has at large \(s\) the form \(c_0+c_1/s+O(s^{-2})\), checked at orbit points with random \(q,\alpha\). (i) \(c_0\) is real and is the Maxwell coefficient of the gauge field \(a_b=\epsilon\zeta\alpha_b\sin\phi\) induced by the kernel momenta: \(c_0=\lambda_0\,G(\alpha)\) with \(G=(\kappa\cdot\kappa)(\alpha\cdot\alpha)-(\kappa\cdot\alpha)^2\) (agreement to \(10^{-8}\) or better in the tensor, \(\lambda_0>0\) at all sampled points). It is gauge invariant (\(\alpha\propto\kappa\)), of the sign of \(G\), which is indefinite in Lorentzian signature, and the integral \(\int^{s_{\max}}\mathrm ds\,c_0\) diverges linearly: the first-order mass term has no finite local coefficient at this order, and the derivative expansion is expected to be non-uniform at small \(\xi\) of the non-kernel directions (a mass insertion into the massless non-kernel propagator; not computed). (ii) \(c_1\) is purely imaginary and, for the window of \(\Delta_{3/4}\) at the sampled points, is neither invariant under \(\alpha\to\alpha+\lambda\kappa\) (relative violation \(0.04\)--\(0.6\)) nor of covariant form, and has eigenvalues of both signs. (iii) Let \(V'\) be the vertical tangent space (common kernel of the \(\mathrm dy^\mu\)) and scale by \(f\) the part of the window quadratic form in \(\omega_a\) in the splitting \(\xi\mapsto(\xi(W_a),\xi(v_k))\). As \(f\to0\), \(c_1\to\mathrm i\lambda_1 A\) with \(\lambda_1=|\det\Lambda|^{-1}(4\pi)^{-3}\det(P'_{kk})^{-1/2}>0\) (\(\Lambda\) the matrix of the splitting, \(P'_{kk}\) the vertical block of the window form) and no free parameter; the difference is \(O(f)\) at four orbit points (relative \(8\cdot10^{-4}\), \(1.4\cdot10^{-3}\), \(7\cdot10^{-4}\), \(8\cdot10^{-5}\) at \(f=10^{-3}\)). Here, with the base cometric \(\bar g^{\mu\nu}=\eta^{ab}W_a(y^\mu)W_b(y^\nu)\), \(h^{\mu\nu}\) its first-order change \(-\sin\phi\,(u^\mu v^\nu+v^\mu u^\nu)\), \(u^\mu=\eta^{ab}\alpha_aW_b(y^\mu)\), \(v^\mu=\{\phi,y^\mu\}\) (so \(q\cdot v=0\)) and \(a_\mu=\bar g_{\mu\nu}u^\nu\),
\[
A=\tfrac1{24}\Bigl[v^2\bigl((q\cdot a)^2-q^2a^2\bigr)-2q^2(a\cdot v)^2\Bigr],
\]
the sum of the Einstein--Hilbert second variation, with the phase-averaged Fierz--Pauli form \(\tfrac16\,\mathrm{FP}[h]=\tfrac1{24}\bigl[v^2\bigl((q\cdot a)^2-q^2a^2\bigr)+q^2(a\cdot v)^2\bigr]\), and the dilaton-gradient term \(-\tfrac18q^2(a\cdot v)^2\) of the Liouville measure. \(A\) is invariant under \(a\to a+\lambda q\) (checked to \(10^{-16}\)), vanishes only on that direction, and is indefinite in general (definite only for \(q\) time-like and \(v\) space-like, where it is positive). (iv) The base two-form \(\theta^{\mu\nu}=\{y^\mu,y^\nu\}=M^{\mu\nu}\) satisfies \(\theta^{\rho\mu}\eta_{\mu\nu}\theta^{\sigma\nu}=c_\theta\eta^{\rho\sigma}+y^\rho y^\sigma-t^\rho t^\sigma\), \(y^\mu=M^{\mu5}\), \(t^\mu=M^{\mu4}\) (to machine precision at six points), with \(c_\theta\) a scalar function of the orbit point, so \(v^2=c_\theta q^2+(q\cdot y)^2-(q\cdot t)^2\) depends on the time-like vector \(y^\mu\) and on the fibre point. After a Haar average over the fibre at a base point with \(y^i=0\), the \(A\)-form (weighted by the Liouville density and \(\lambda_1\)) remains gauge invariant but is not proportional to \(q^2\eta_{ab}-q_aq_b\) (remainder \(0.8\)--\(1\) of its norm), and its covariant part changes by a factor \(2\) between boosts of rapidity \(0\) and \(0.8\) at fixed \(q\cdot q\). The window and the leaf both single out \(y^\mu\), which this check does not separate.

\emph{What is not shown.} The Pauli--Sakharov completion (the term \(M_{\mathrm{str}}\) multiplying the logarithm of the cutoff is accompanied by \(-\sum_iw_id_im_i\ln m_i\) when the masses are kept non-perturbatively) is expected from the Euclidean heat-kernel structure and is not derived for this hypoelliptic operator. The background flux, which cuts off large \(s\) in the full leaf, is not included, so the divergence in (i) is not resolved.

This is a statement for the Poisson limit, one loop, the term linear in the masses, the order \(\epsilon^2q^4\) with the first jet frozen, the supplied content with no bilinear fermion mass, the Gaussian window of \(\Delta_{3/4}\), and numerical sampling; the limit in (iii) is a limit of the window family and not the window of the class. It identifies no induced form with the emergent metric of Chapter~\ref{ch:vierbein-emergence}, makes no Wick-rotation claim, and gives no value or sign of \(G_{\mathrm{eff}}\). This changes no ledger label; Problem~F remains open.
